% Generated by tools/edn_latex.py from reports/edn.md.
% Do not edit: change reports/edn.md and run `make edn`.
\ednentry{
  id = {EDN-02},
  title = {Model family: gradient boosting (LightGBM as main model)},
  date = {2026-09-22},
  milestone = {M1: Project Inception},
  topic = {Modelling Approach},
  participants = {All team members},
  decision = {Use gradient boosting as the model family, with LightGBM as the main model (basic features, then all features) and CatBoost as a challenger for high-cardinality categoricals.},
  rationale = {\emph{Alternatives considered:}
\begin{itemize}[leftmargin=*, topsep=2pt]
\item \textbf{Option A (chosen): Gradient boosting (LightGBM main, CatBoost challenger).}
\newline Pros: best-in-class on medium tabular data, native categoricals and missing values, trains in minutes on CPU, small artefacts, exact SHAP explanations.
\newline Cons: more tuning surface than a linear baseline; still needs a simpler baseline (median, Ridge) to prove its value.
\item \textbf{Option B: Linear / Ridge regression only.}
\newline Pros: simple, fast, highly interpretable depreciation baseline.
\newline Cons: cannot capture non-linear interactions (brand x age x mileage) well, lower expected accuracy.
\item \textbf{Option C: kNN, random forest, tabular NNs, hierarchical Bayes, LLM zero-shot.}
\newline Pros: each offers some property (kNN gives comparables for free, NNs could use text embeddings, LLM needs no training).
\newline Cons: kNN and random forest scale poorly or lag gradient boosting on this data size; tabular NNs need more data/tuning and lose native categorical handling; hierarchical Bayes is slow to fit at this scale; LLM zero-shot has no calibrated uncertainty and conflicts with the CPU-only, small-image constraint.
\end{itemize}
\par
\emph{Why:} The course grades MLOps practice over raw accuracy, so the team wants a model that trains fast on CPU, handles categoricals/missing values natively, keeps artefacts small, and supports exact SHAP explanations for UC1/UC2 -- gradient boosting (LightGBM) is the best fit, with a median and Ridge baseline kept in the experiment ladder to prove it earns the added complexity.},
  ai = {Information seeking, Alternative assessment},
  response = {Accepted},
  assessment = {AI laid out the alternatives and the CPU/size/explainability constraints behind the ``why gradient boosting'' rationale in the project brief; the team reviewed and confirmed the choice, including keeping CatBoost as a challenger rather than the main model, in today's sprint planning.},
  aievidence = {Modelling plan and ``Why gradient boosting'' write-up in \texttt{docs/\allowbreak{}docs/\allowbreak{}project-\allowbreak{}brief.\allowbreak{}md} \S{}4 (2026-09-22), confirmed in today's sprint planning session.},
  evidence = {\href{https://github.com/mlops-2627q1-mds-upc/MLOPS_Recommenditos/pull/13}{PR \#13}.},
}
